\chapter{Solutions to Chapter 9: Bridging Finite and Super Population Causal Inference}
\label{sol:chapter09}

\begin{exercisesolution}{OLS decomposition of the observed outcome under the CRE}{由两个 potential-outcome projections 合成 observed-outcome projection}
本题的核心是一个简单但很有用的代数事实：在 CRE 的 superpopulation 视角下，observed outcome
\[
Y=ZY(1)+(1-Z)Y(0)
\]
把两个 potential outcomes 的线性投影拼接在同一个 fully interacted regression 里。由 \cref{eq::ols-decomposition-y1,eq::ols-decomposition-y0}，
\[
Y(1)=\gamma_1+\beta_1\tran X+\varepsilon(1),
\qquad
Y(0)=\gamma_0+\beta_0\tran X+\varepsilon(0),
\]
其中 OLS projection 的 normal equations 给出
\[
\mathbb{E}\varepsilon(1)=0,\quad
\mathbb{E}\{X\varepsilon(1)\}=0,\qquad
\mathbb{E}\varepsilon(0)=0,\quad
\mathbb{E}\{X\varepsilon(0)\}=0.
\]
把这两个分解代入 observed outcome，可得
\begin{eqnarray*}
Y
&=&Z\{\gamma_1+\beta_1\tran X+\varepsilon(1)\}
 +(1-Z)\{\gamma_0+\beta_0\tran X+\varepsilon(0)\}\\
&=&\gamma_0+(\gamma_1-\gamma_0)Z
  +\beta_0\tran X+(\beta_1-\beta_0)\tran XZ
  +Z\varepsilon(1)+(1-Z)\varepsilon(0).
\end{eqnarray*}
因此只要定义
\[
\alpha_0=\gamma_0,\qquad
\alpha_Z=\gamma_1-\gamma_0,\qquad
\alpha_X=\beta_0,\qquad
\alpha_{ZX}=\beta_1-\beta_0,
\]
以及
\[
\varepsilon=Z\varepsilon(1)+(1-Z)\varepsilon(0),
\]
就得到题目中的 observed-outcome decomposition
\[
Y=\alpha_0+\alpha_ZZ+\alpha_X\tran X+\alpha_{ZX}\tran XZ+\varepsilon.
\]

剩下要说明的是：这不仅是一个代数恒等式，而且确实是 population OLS projection。记 regressors 为
\[
W=(1,\ Z,\ X\tran,\ ZX\tran)\tran.
\]
OLS projection 的判别条件是 normal equations
\[
\mathbb{E}(W\varepsilon)=0.
\]
在 \cref{def::cre-superpopulation} 的 CRE 条件下，$Z\ind\{Y(1),Y(0),X\}$，从而 $Z$ 与 $\varepsilon(1),\varepsilon(0),X$ 独立。于是
\[
\mathbb{E}\varepsilon
=\mathbb{E}\{Z\varepsilon(1)\}
 +\mathbb{E}\{(1-Z)\varepsilon(0)\}
=\mathbb{E}Z\,\mathbb{E}\varepsilon(1)
 +\mathbb{E}(1-Z)\,\mathbb{E}\varepsilon(0)
=0.
\]
同理，
\[
\mathbb{E}(Z\varepsilon)
=\mathbb{E}\{Z\varepsilon(1)\}
=\mathbb{E}Z\,\mathbb{E}\varepsilon(1)
=0,
\]
因为 $Z(1-Z)=0$。对 covariate regressors，也有
\begin{eqnarray*}
\mathbb{E}(X\varepsilon)
&=&\mathbb{E}\{ZX\varepsilon(1)\}
  +\mathbb{E}\{(1-Z)X\varepsilon(0)\}\\
&=&\mathbb{E}Z\,\mathbb{E}\{X\varepsilon(1)\}
  +\mathbb{E}(1-Z)\,\mathbb{E}\{X\varepsilon(0)\}
=0,
\end{eqnarray*}
以及
\[
\mathbb{E}(ZX\varepsilon)
=\mathbb{E}\{ZX\varepsilon(1)\}
=\mathbb{E}Z\,\mathbb{E}\{X\varepsilon(1)\}
=0.
\]
这些正交条件说明
\[
(\alpha_0,\alpha_Z,\alpha_X,\alpha_{ZX})
=\arg\min_{a_0,a_Z,a_X,a_{ZX}}
\mathbb{E}\{Y-a_0-a_ZZ-a_X\tran X-a_{ZX}\tran XZ\}^2,
\]
即题目所要求的 OLS decomposition。

\paragraph{Remark.}
这个练习解释了为什么 fully interacted OLS 是 Lin estimator 的自然 population analog。若先分别投影 $Y(1)$ 和 $Y(0)$，再把 treatment indicator 放回 observed outcome，interaction coefficient 必然等于两组 projection slopes 的差 $\beta_1-\beta_0$。因此，在 \cref{sec::cre-superpop} 中，当 covariates 被中心化时，pooled interacted regression 中 $Z$ 的系数正好估计
\[
\gamma_1-\gamma_0,
\]
也就是 \citet{lin2013} 的 estimator 所对应的 population quantity。
\end{exercisesolution}

\begin{exercisesolution}{Variance estimation of \citet{lin2013}'s estimator under the super population framework}{四种 standard errors 的 simulation 比较}
这道题的目的不是证明一个新的闭式公式，而是用 simulation 看清楚 \cref{eq::lin-ehw-correction} 的来源。Finite population analysis 中，Lin estimator 的 EHW robust variance 主要反映 treatment assignment 的随机性；superpopulation analysis 中，covariate sample mean 本身也在随机波动。当 $\beta_1\neq\beta_0$ 时，这个额外波动的一阶贡献正是
\[
(\beta_1-\beta_0)\tran\text{var}(X)(\beta_1-\beta_0)/n,
\]
样本版本就是 \cref{eq::lin-ehw-correction}。

下面给出一个完整可复现的 simulation design。取 two-dimensional covariate $X=(X_1,X_2)\tran$，令
\[
Y(0)=1+X_1+0.5X_2+u_0,
\]
\[
Y(1)=1+\tau+2X_1-X_2+u_1.
\]
于是
\[
\beta_0=(1,\ 0.5)\tran,\qquad
\beta_1=(2,\ -1)\tran,
\]
所以 $\beta_1\neq\beta_0$。为简单起见，设 $\tau=0$，且 $u_0,u_1$ 独立于 $X,Z$。CRE under superpopulation framework 使用 \cref{def::cre-superpopulation}：
\[
Z\ind\{Y(1),Y(0),X\}.
\]

对每个 Monte Carlo dataset，计算四个 variance estimates：
\begin{enumerate}
  \item $V_{\mathrm{EHW}}$：fully interacted OLS 中 $Z$ 系数的 EHW robust variance；
  \item $V_{\mathrm{corr}}$：$V_{\mathrm{EHW}}$ 加上 \cref{eq::lin-ehw-correction}；
  \item $V_{\mathrm{boot,fix}}$：bootstrap variance，其中 covariates 只在原始 dataset 中中心化一次，bootstrap sample 不重新中心化；
  \item $V_{\mathrm{boot,recenter}}$：bootstrap variance，其中每个 bootstrap sample 都用该 sample 自身的 covariate mean 重新中心化。
\end{enumerate}
True variance 用 outer Monte Carlo repetitions 中 $\hat\tau_{\textsc{L}}$ 的 empirical variance 估计。

可运行的 \ri{R} 代码如下。
\begin{lstlisting}
library(sandwich)

make_data <- function(n = 500, pz = 0.5) {
  X <- matrix(rnorm(n * 2), n, 2)
  u0 <- rnorm(n, sd = 1)
  u1 <- rnorm(n, sd = 1)
  tau <- 0

  Y0 <- 1 + X[, 1] + 0.5 * X[, 2] + u0
  Y1 <- 1 + tau + 2 * X[, 1] - X[, 2] + u1
  Z <- rbinom(n, 1, pz)
  Y <- Z * Y1 + (1 - Z) * Y0

  data.frame(Y = Y, Z = Z, X1 = X[, 1], X2 = X[, 2])
}

lin_fit <- function(dat, recenter = TRUE) {
  X <- as.matrix(dat[, c("X1", "X2")])
  if (recenter) {
    X <- scale(X, center = TRUE, scale = FALSE)
  }
  fitdat <- data.frame(Y = dat$Y, Z = dat$Z,
                       X1 = X[, 1], X2 = X[, 2])
  fit <- lm(Y ~ Z * (X1 + X2), data = fitdat)
  est <- coef(fit)["Z"]
  vehw <- vcovHC(fit, type = "HC2")["Z", "Z"]
  inter <- coef(fit)[c("Z:X1", "Z:X2")]
  vcorr <- vehw + drop(t(inter) %*% cov(X) %*% inter) / nrow(dat)
  c(est = est, vehw = vehw, vcorr = vcorr)
}

boot_var <- function(dat, B = 400, recenter_boot = TRUE) {
  n <- nrow(dat)
  X0 <- scale(as.matrix(dat[, c("X1", "X2")]),
              center = TRUE, scale = FALSE)
  dat0 <- data.frame(Y = dat$Y, Z = dat$Z,
                     X1 = X0[, 1], X2 = X0[, 2])

  vals <- replicate(B, {
    id <- sample.int(n, n, replace = TRUE)
    if (recenter_boot) {
      lin_fit(dat[id, ], recenter = TRUE)["est"]
    } else {
      lin_fit(dat0[id, ], recenter = FALSE)["est"]
    }
  })
  var(vals)
}

one_run <- function(n = 500, Bboot = 400) {
  dat <- make_data(n)
  lf <- lin_fit(dat, recenter = TRUE)
  c(est = unname(lf["est"]),
    vehw = unname(lf["vehw"]),
    vcorr = unname(lf["vcorr"]),
    vboot_fix = boot_var(dat, B = Bboot, recenter_boot = FALSE),
    vboot_recenter = boot_var(dat, B = Bboot, recenter_boot = TRUE))
}

set.seed(20260606)
M <- 1000
ans <- replicate(M, one_run(n = 500, Bboot = 400))

true_var <- var(ans["est", ])
c(true_var = true_var,
  mean_ehw = mean(ans["vehw", ]),
  mean_corrected = mean(ans["vcorr", ]),
  mean_boot_fix = mean(ans["vboot_fix", ]),
  mean_boot_recenter = mean(ans["vboot_recenter", ]))

cover <- c(
  ehw = mean(abs(ans["est", ]) <=
             1.96 * sqrt(ans["vehw", ])),
  corrected = mean(abs(ans["est", ]) <=
                   1.96 * sqrt(ans["vcorr", ])),
  boot_fix = mean(abs(ans["est", ]) <=
                  1.96 * sqrt(ans["vboot_fix", ])),
  boot_recenter = mean(abs(ans["est", ]) <=
                       1.96 * sqrt(ans["vboot_recenter", ]))
)
cover
\end{lstlisting}

这段 simulation 应呈现如下模式。第一，$V_{\mathrm{EHW}}$ 的平均值通常小于 \ri{true\_var}，因为它没有计入 random covariate mean 带来的 superpopulation variation。第二，$V_{\mathrm{corr}}$ 的平均值应接近 \ri{true\_var}；这正是 \cref{eq::lin-ehw-correction} 的作用。第三，若 bootstrap sample 不重新中心化 covariates，它在每个 bootstrap world 中把原始样本的 covariate center 当作固定对象，容易漏掉一部分 centering uncertainty。第四，每个 bootstrap sample 重新中心化时，bootstrap procedure 更忠实地模拟了实际估计步骤，因此其平均 variance 通常更接近 corrected variance 和 true variance。

\paragraph{Remark.}
这道题把 finite population 和 superpopulation 的差异压缩到一个可观察的现象：同一个 Lin regression，在 fixed covariates 的设计随机化分析中可以用 EHW robust variance；但当 covariates 也是从 superpopulation 抽出来时，$\bar X$ 的随机性不能忽略。只要 $\beta_1=\beta_0$，这个额外项消失；一旦 $\beta_1\neq\beta_0$，忽略它就会低估 uncertainty。
\end{exercisesolution}
